\chapter{The sample and its law}
\label{ch:sample}

\begin{definition}[The sample of an offspring field]
\label{def:sample}
\uses{def:word, def:subtree}
\lean{BranchingProcess.sample, BranchingProcess.Survives}
\leanok
A field of offspring counts is a map $c\colon\cN(N)\to\bbN$. The \emph{sample} of $c$ is the
subtree in which each vertex $v$ keeps its first $c(v)$ children,
\[
  \cT(c)=\set{v : \text{every letter of } v \text{ is below the count at the vertex it leaves}}.
\]
The field \emph{survives} when $\cT(c)$ is infinite.
\end{definition}

\begin{lemma}[The branching property]
\label{lem:branching}
\uses{def:sample}
\lean{BranchingProcess.append_mem_sample_iff, BranchingProcess.subAt_sample,
BranchingProcess.coe_sample}
\leanok
The residual subtree of $\cT(c)$ at a vertex of the sample is the sample of the shifted
field, and $\cT(c)$ decomposes into its root together with the samples hanging off the
children of the root.
\end{lemma}

\begin{proof}
\leanok
Both are read off the definition: membership of $vw$ imposes the conditions of $v$ followed by
those of $w$ for the shifted field.
\end{proof}

\begin{definition}[The skeleton]
\label{def:skeleton}
\uses{def:sample}
\lean{BranchingProcess.skeleton}
\leanok
The \emph{skeleton} of a sample is the set of its vertices whose residual subtree is infinite.
\end{definition}

\begin{lemma}[Rays in the skeleton]
\label{lem:skeleton-ray}
\uses{def:skeleton, lem:treeDist}
\lean{BranchingProcess.exists_child_mem_skeleton, BranchingProcess.exists_skeleton_ray,
BranchingProcess.exists_isometric_skeleton_ray}
\leanok
A skeleton vertex has a child in the skeleton, and hence carries a ray, an isometric copy of
$\bbN$ inside the sample.
\end{lemma}

\begin{proof}
\leanok
A vertex with infinite residual subtree has, by the pigeonhole principle over its finitely
many children, a child with infinite residual subtree. Iterating produces the ray, and
distances along it are computed by \cref{lem:treeDist}.
\end{proof}

\begin{definition}[The law of the sample]
\label{def:sampleMeasure}
\uses{def:offspring, def:sample}
\lean{BranchingProcess.fieldMeasure, BranchingProcess.sampleMeasure, BranchingProcess.treeLaw}
\leanok
The \emph{field measure} on $\cN(N)\to\bbN$ is the infinite product of copies of $\theta$, one
per vertex. Its pushforward along $\cT$ is the \emph{law of the Galton--Watson tree} of
$\theta$, a measure on subtrees for the $\sigma$-algebra generated by the vertex events
$\set{T : v\in T}$.
\end{definition}

\begin{theorem}[The extinction probability is the probability of extinction]
\label{thm:extinction}
\uses{def:sampleMeasure, def:extinction, lem:branching}
\lean{BranchingProcess.sampleMeasure_not_survives, BranchingProcess.sampleMeasure_survives}
\leanok
If $J\leq N$, then $\bbP(\cT \text{ finite})=q$.
\end{theorem}

\begin{proof}
\leanok
Let $q_n$ be the probability that generation $n$ is empty. The event that generation $n+1$ is
empty splits at the root into the offspring count there and $N$ independent copies of the
process, so $q_{n+1}=f(q_n)$ with $q_0=0$. The iterates $f^n(0)$ increase to the least fixed
point of $f$ in $[0,1]$, which is $q$ by \cref{lem:extinction-fixed}. Extinction is the
increasing union of the events that some generation is empty.
\end{proof}
